\section{Preliminaries}
\label{sec:basic}

We begin by introducing the syntax and the semantics of $\Hxpath$.
We assume a fixed \emph{signature} of pairwise disjoint sets: $\props = \set{p, q, r, \dots}$ of \emph{propositional} symbols, $\noms = \set{i, j, k, \dots}$ of \emph{nominals}, $\rels = \set{\astep[a], \astep[b], \astep[c], \dots }$ of symbols for \emph{accessibility relations}, and $\cmps = \set{\cmp, \mathsf{f}, \mathsf{g},\dots}$ of \emph{comparisons by data (in)equality}.
We further assume that $\props$ and $\noms$ are countably infinite, and that $\rels$ and $\cmps$ are finite.
    
\begin{defi}\label{def:language}
  The \emph{language} of $\Hxpath$ consists of \emph{path} expressions $\alpha$, $\beta$, $\gamma$, \dots, and \emph{node} expressions $\varphi$, $\psi$, $\chi$, \dots.
  Path and node expressions are defined by mutual recursion according to the grammar:
  \begin{align*}
    \alpha, \beta & \Coloneqq
        \pathh \mid
        \goto{i} \mid
        \varphi? \mid
        \alpha \beta \mid
        \alpha \cup \beta \\
    \varphi, \psi & \Coloneqq
        p \mid
        i \mid
        \lnot \varphi \mid
        \varphi \land \psi \mid
        i{:}\varphi \mid 
        \tup{\child} \varphi \mid
        \tup{\alpha =_{\cmp} \beta} \mid
        \tup{\alpha \neq_{\cmp} \beta},
  \end{align*}
  \noindent where $p\in\props$, $i\in\noms$, $\pathh\in\rels$, and $\cmp\in\cmps$.
  By way of notation, we use~$\cmpr$ when there is no need to distinguish between $=_\cmp$ and $\neq_\cmp$.
  We sometimes refer to node expressions of the form $i{:}\varphi$, $\tup{@_i\alpha \cmpr @_j\beta}$, or $\lnot\tup{@_i\alpha \cmpr @_j\beta}$, as \emph{prefixed expressions}. 
\end{defi}

% Notice that `$i{:}$' is overloaded.
% By itself is a path expression; when it is followed by a path expression $\alpha$, the result $i{:}\alpha$ is a path expression; when it is followed by a node expression $\varphi$, the result $i{:}\varphi$ is a node expression.
% In any case, its use will be clear from the context.
%, and $\cmpr$ when the subindex $\cmp$ is irrelevant.

In the syntax above we find nominals and satisfiability statements. The latter come in two flavours: as path expressions of the form $@_i\alpha$, and as node expressions of the form $i{:}\varphi$. The modal diamond $\tup{\child}\varphi$ is included as primitive (as we will see, this has a positive impact in the definition of tableaux in~\Cref{sec:tableaux}).  Notice that we are in a multi-modal system, as we have one modal diamond for each $\pathh\in\rels$. We also allow for different (in)equality comparisons, one for each $\cmp\in\cmps$. 
%
Finally, we make use the following abbreviations for node expressions:
\begin{align*}
  \varphi \to \psi & \coloneqq \lnot\varphi \lor \psi   &
  \varphi \liff \psi & \coloneqq (\varphi \to \psi) \land (\psi \to \varphi) &
  \varphi \lor \psi & \coloneqq \lnot(\lnot\varphi \land \lnot\psi)  \\
  \tup{\alpha}\varphi & \coloneqq \tup{\alpha\varphi? = \alpha\varphi?} &
  [\alpha]\varphi & \coloneqq \lnot\tup{\alpha}\lnot\varphi &
  [\alpha \cmpr \beta] & \coloneqq \lnot\tup{\alpha \cmpd \beta}.
\end{align*}
The symbol ${\cmpd}$ above is ${\neq_\cmp}$ if ${\cmpr}$ is ${=_\cmp}$, and is ${=_\cmp}$ if ${\cmpr}$ is ${\neq_\cmp}$.
%Finally, we use $\epsilon$ as an abbreviation for the path expression $\top?$.

% Following the standard notation in XPath logics and in modal logics, the $[\ ]$ operation is overloaded: for  $\pathh\in\rels$, $\varphi$ a node expression and $\alpha$ a path expression, both $[\pathh]\varphi$ and $[\varphi]\pathh$ are well-formed expressions; the former is a node expression where $[\pathh]$ is a box modality, the latter is a path expression where $[\varphi]$ is a test.

The semantics of path and node expressions is given with respect to \emph{hybrid data models}.

\begin{defi}\label{def:models}
  A \emph{(hybrid data) model} is a tuple $\model = \tup{N, \set{R_\pathh}_{\pathh \in \rels}, \set{\approx_{\cmp}}_{\cmp\in\cmps}, \valuation, \naming}$ where:
  \begin{enumerate} 
    \item $N$ is a non-empty set of elements called \emph{nodes}, 
    \item each $R_{\pathh}\subseteq {N\times N}$ is a binary \emph{accessibility} relation,
    \item each ${\approx_{\cmp}} \subseteq {N\times N}$ is an \emph{equivalence} relation, 
    \item $\valuation :\props \to 2^N$  is a \emph{valuation} function, and 
    \item $\naming : \noms \to N$ is a \emph{naming} function.
  \end{enumerate}
\end{defi}

In brief, the models of $\Hxpath$ can be understood as the typical models of Hybrid Modal Logic, see, e.g.,~\cite{arec:hybr05b}, augmented with a collection of equivalence relations intended to capture equality criteria between data values in nodes.
The \emph{satisfiability} relation is defined below. Notice how satisfiability of a path expression on a model is relative to two nodes (as it describes the accessibility from one to the other), whereas satisfiability of a node expression is at a particular node (stating truth at that node). 
%Before moving to the formal definition, let us provide a brief explanation of the intuitive semantics of path and node expressions.
%We begin by assuming an underlying graph structure with nodes and edges.
%Then, p

%Path expressions, as their name indicates, are meant to represent directed ``paths'' from a starting node to an end node in a graph.
Precisely, the path expression $\pathh$ can be understood as traversing an edge labelled by $\pathh$ in the graph.
In turn, $\goto{i}$ can be understood as a ``go to'' statement, i.e., jump to \emph{the} node named by the nominal~$i$.
Then, $\varphi?$ can be understood as a ``test on the truth of $\varphi$''. % akin to those in Propositional Dynamic Logic.
Finally, $\alpha\beta$ can be understood as the ``concatenation'' of $\alpha$ and $\beta$, 
%i.e., follow $\alpha$ and then follow $\beta$, 
and $\alpha \cup \beta$ as the ``non-deterministic choice'' between $\alpha$ and $\beta$. %, i.e., follow $\alpha$ or follow $\beta$. 
The intuitive semantics of the node expressions of the form $p$, $i$, $\lnot\varphi$, $\varphi \land \psi$, $i{:}\varphi$, and $\tup{\pathh}\varphi$ are the same as in Hybrid Modal Logic (see, e.g.,~\cite{arec:hybr05b} for details). 
%In this case, the intuitive semantics of $p$, $\lnot \varphi$, and $\varphi \land \psi$ is as expected when thought of at a particular node in the graph.
Specifically, %$i$ can be understood as a node's name, whereas 
$i{:}\varphi$ can be understood as ``$\varphi$ is true at the node named by $i$'', while $\tup{\pathh}\varphi$ is read as ``$\varphi$ is true after traversing some  $\pathh$-edge''.
Finally, the intuitive semantics of data equalities $\tup{\alpha =_{\cmp} \beta}$ (and data inequalities $\tup{\alpha \neq_{\cmp} \beta}$) is ``there are $\alpha$ and $\beta$ paths with a common starting point whose end points have data values that are (not) equal according to the criteria determined by $\cmp$''.

\begin{defi}\label{def:semantics}
  Let $\model = \tup{N, \set{R_\pathh}_{\pathh \in \rels}, \set{\approx_{\cmp}}_{\cmp\in\cmps}, \valuation, \naming}$ be a model, and $n,n'\in N$. Then: 
  \[
    \begin{array}{l@{\quad}l@{\quad}l}
      \model, n, n' \vDash \child
        & \text{\it iff} &   n R_{\child} n'   \\
      \model, n, n' \vDash \goto{i}
        & \text{\it iff } &  \naming(i) = n' \\ 
      \model, n, n' \vDash  \varphi?
        & \text{\it iff} & n=n' \text{ and } \model, n \vDash   \varphi  \\
      \model, n, n' \vDash \alpha\beta
        & \text{\it iff} & \text{exists } z \text{ s.t.\ } \model,  n, z \vDash \alpha \text{ and }  \model,  z, n' \vDash \beta \\
      \model, n, n' \vDash \alpha \cup \beta
        & \text{\it iff} & \model, n, n' \vDash  \alpha  \text{ or } \model, n, n' \vDash  \beta \\
      \model, n    \vDash p
        & \text{\it iff} &  n \in \valuation(p)\\
      \model, n    \vDash i
        & \text{\it iff} &  \naming(i) = n \\
      \model, n    \vDash \lnot \varphi
        & \text{\it iff} &   \model, n \nvDash  \varphi \\
      \model, n    \vDash \varphi \land \psi
        & \text{\it iff} &   \model, n \vDash  \varphi \text{ and }  \model, n \vDash  \psi  \\
      \model, n    \vDash  \tup{\pathh}\varphi
        & \text{\it iff} & \text{exists } z \text{ s.t.\ } nR_{\pathh}z \text{ and } \model, z \vDash \varphi \\
      \model, n   \vDash  i{:}\varphi
        & \text{\it iff } & \model, \naming(i) \vDash \varphi \\
      \model, n    \vDash \tup{\alpha =_\cmp \beta}
        & \text{\it iff} & \text{exist } z, z' \text{ s.t.\ } \model, n, z \vDash  \alpha \text{, } \model, n, z' \vDash  \beta \text{, and } z \approx_\cmp z' \\
      \model, n    \vDash \tup{\alpha \neq_\cmp \beta}
        & \text{\it iff} & \text{exist } z, z' \text{ s.t.\ } \model, n, z \vDash  \alpha \text{, } \model, n, z' \vDash  \beta \text{, and } z \not\approx_\cmp z'.
    \end{array}
  \]
  A node expression $\varphi$ is satisfiable iff there is $\model, n$ s.t., $\model, n \vDash \varphi$.
  %A set $\Sigma$ of node expressions is \emph{satisfiable} iff for all $\varphi \in \Sigma$, $\varphi$ is satisfiable.
  % We write $\model\vDash\varphi$, and say that $\varphi$ is \emph{valid} in $\model$, iff $\model,n\vDash\varphi$, for every node $n$ of $\model$.
  In turn, $\varphi$ is \emph{valid}, written $\vDash \varphi$ iff $\model, n \vDash\varphi$, for all $\model, n$.
  These notions are analogously defined for path expressions. 
  %We say that $\varphi$ is a \emph{semantic consequence} of the set of premisses $\Sigma$, and write $\Sigma \vDash \varphi$, iff $\Sigma \cup \set{\lnot \varphi}$ is not satisfiable.
  %We write $\vDash \varphi$, and call $\varphi$ a \emph{tautology}, if $\emptyset \vDash \varphi$.
\end{defi}



We conclude this section by listing some immediate properties for node expressions.
These properties are useful in the justification of tableaux rules in \Cref{def:tableau}.

\begin{prop}\label{prop:equivalences}
  Let $\model$ be any model. The following node expressions are valid in $\model$:
  \begin{align*}
    i{:}\tup{\alpha \cmpr \beta}
      & \liff \tup{@_i\alpha \cmpr @_i\beta} &
    i{:}\lnot\tup{\alpha \cmpr \beta}
      & \liff \lnot\tup{@_i\alpha \cmpr @_i\beta} \\
    \tup{\varphi?\alpha \cmpr \beta}
      & \liff \varphi \land \tup{\alpha \cmpr \beta} &
    \lnot\tup{\varphi?\alpha \cmpr \beta}
      & \liff \lnot\varphi \lor \lnot\tup{\alpha \cmpr \beta} \\
    \tup{(\alpha \cup \beta)\gamma \cmpr \eta}
      & \liff \tup{\alpha\gamma \cmpr \eta} \lor \tup{\beta\gamma \cmpr \eta} &
    \lnot\tup{(\alpha \cup \beta)\gamma \cmpr \eta} & \liff \lnot\tup{\alpha\gamma \cmpr \eta} \land \lnot\tup{\beta\gamma \cmpr \eta}.
  \end{align*}
\end{prop}